% =====================================================================
%  2.4 Ulikheter — ekstra utfordring
% =====================================================================

\begin{infoboks}{Velg metode selv}
Arket bygger på 2.4 Ulikheter. Bruk fortegnslinje, graf eller regning, og velg selv.
Vis regning og begrunnelse. Kontroller svaret ved å sette inn tall.
\end{infoboks}

\nivaahode{2}

\oppg{U1}{2}{\uh}
Løs $x^3 - 4x \ge 0$. Fortegnslinja får tre faktorer.

\oppg{U2}{2}{\mh}
Høyden til en ball er $h(t) = 20t - 5t^2$ meter, der $t$ er tiden i sekunder etter kastet.
\begin{tasks}(1)
  \task I hvilket tidsrom er ballen høyere enn 15 m?
  \task Hvor lenge er ballen høyere enn 15 m?
\end{tasks}

\oppg{U3}{2}{\prog}
Programmet tester om $x^2 - 2x - 3 < 0$ for noen verdier av $x$.
\begin{kode}
x = -10
while x <= 10:
    if x**2 - 2*x - 3 < 0:
        print(x)
    x = x + 0.5
\end{kode}
\begin{tasks}(1)
  \task Hva skriver programmet ut?
  \task Hvorfor gir utskriften ikke hele løsningen? Lag en ulikhet der programmet går glipp av løsninger.
  \task Endre $<$ til $\le$ i programmet. Hvilke verdier kommer i tillegg, og hva forteller det om endepunktene?
\end{tasks}

\oppg{U4}{2}{\uh}
En elev løser $x^2 < 3x$ slik: «Jeg deler på $x$ og får $x < 3$.»
\begin{tasks}(1)
  \task Forklar hva som er galt. Bruk et tall som viser det.
  \task Løs ulikheten riktig.
\end{tasks}

\nivaahode{3}

\oppg{U5}{3}{\uh}
Løs $\dfrac{x - 1}{x + 2} \ge 0$ med fortegnslinje. Teller og nevner får hver sin rad.
Forklar hvorfor $x = -2$ aldri er med i løsningen, mens $x = 1$ er med.

\oppg{U6}{3}{\uh}
Anta at $a, b \ge 0$.
\begin{tasks}(1)
  \task Vis at $\dfrac{a + b}{2} \ge \sqrt{ab}$. Start med at $(\sqrt{a} - \sqrt{b})^2 \ge 0$.
  \task Når er det likhet? Begrunn.
  \task Av alle rektangler med omkrets 20 cm: Hvilket har størst areal? Bruk resultatet i a).
\end{tasks}

\oppg{U7}{3}{\uh}
Undersøk ulikheten $x^2 - 4x + k \le 0$ for ulike verdier av $k$.
For hvilke $k$ har den minst én løsning? Beskriv løsningsmengden i hvert tilfelle, og begrunn med diskriminanten.

\begin{oppsummering}
Samle alt på én side og faktoriser. Bruk fortegnslinje, og del aldri på noe som kan være negativt.
Diskriminanten sier om parabelen krysser $x$-aksen.
Kontroller svaret ved å sette inn et tall fra hvert område.
\end{oppsummering}
